\chapter{Performance Measurement}\label{ch:rs:performance-measurement}

Two strategies have Sharpe ratios of 1.00 and 1.09 over five years of daily returns. The first lost 33.6 per cent peak to trough and spent 791 trading
days, more than three years, under water. The second sells out-of-the-money puts and never had a losing quarter; its worst drawdown was 14.8 per cent.
Both histories were drawn from the chapter's simulations, of strategies with long-run volatilities of 15 and 17 per cent, and the second was chosen
among the 9\% of its simulated five-year histories without a losing quarter. Over the next five years the first has a 14\% chance of a drawdown of 30
per cent or more, the second 41\%. The Sharpe ratio is the right first number and a dangerous last one. This chapter builds the rest of the tear
sheet (standard errors, drawdowns, turnover, the shape of outcomes, the benchmark) as \texttt{firm.perf}, which reads the \texttt{BacktestResult} of
chapters 16 to 18, and reads it for four synthetic streams and chapter~16's momentum book.

\section{The Sharpe ratio and its standard error}

The Sharpe ratio (Book~4, chapter~11) is the mean excess return over its standard deviation; Sharpe introduced it in 1966 as the reward-to-variability
ratio and gave it its present name in 1994.

\begin{definition}[Annualisation]\label{def:rs:performance-measurement:annualisation}
\emph{Annualisation}\index{annualisation} converts a statistic measured on periods of one length to a year: with $q$ periods a year, the mean is multiplied by $q$,
the volatility by $\sqrt q$ and the Sharpe ratio by $\sqrt q$. The scaling of the volatility and of the Sharpe ratio holds only when the returns are
serially uncorrelated.
\end{definition}

The standard error of an estimated Sharpe ratio, for independent normal returns, is about $\sqrt{(1 + \mathrm{SR}^2/2)/T}$ per period over $T$ periods (Lo,
2002; Book~4, chapter~11). Annualised, this is nearly $\sqrt{q/T}$: the length of the history in years sets it, not the frequency. Five years of daily returns
give a standard error of 0.45, whatever the strategy; a measured Sharpe ratio of 1.0 has a 95\% interval from 0.12 to 1.88. The four five-year streams of the chapter
have this iid standard error (the tear sheet in the tutorial); the HAC standard error, which also allows for serial correlation and fat tails, is close for the
trend and short-volatility streams, larger for the market maker's fat-tailed days (0.52), and largest for a smoothed book.

\begin{definition}[Return smoothing, autocorrelation-adjusted Sharpe ratio]\label{def:rs:performance-measurement:smoothing}
\emph{Return smoothing}\index{return smoothing} is the reporting of a moving average of true returns, $r^o_t = \sum_{j=0}^{k}\theta_j r_{t-j}$ with
$\theta_j \ge 0$ and $\sum_j\theta_j = 1$, as happens when illiquid assets are marked at stale or appraised prices (Getmansky, Lo and Makarov). The
\emph{autocorrelation-adjusted Sharpe ratio}\index{autocorrelation-adjusted Sharpe ratio} annualises a per-period Sharpe ratio $\mathrm{SR}_1$ with the
autocorrelations $\rho_k$ of the returns: $\mathrm{SR}(q) = \mathrm{SR}_1\, q\big/\sqrt{q + 2\sum_{k=1}^{q-1}(q-k)\rho_k}$ (Lo, 2002).
\end{definition}

\begin{proposition}[What smoothing does]\label{prop:rs:performance-measurement:smoothing}
If true returns are independent with mean $\mu$ and variance $\sigma^2$, reported returns smoothed with profile $\theta$ have mean $\mu$, variance
$\sigma^2\sum_j\theta_j^2$ and autocorrelations $\rho_k = \sum_j\theta_j\theta_{j+k}\big/\sum_j\theta_j^2$. The naive Sharpe ratio is inflated by
$1/\sqrt{\sum_j\theta_j^2}$; the autocorrelation-adjusted ratio, computed with the population autocorrelations, removes the inflation up to terms of relative
order $1/q$.
\end{proposition}

\begin{proof}
The mean and the autocovariances $\sigma^2\sum_j\theta_j\theta_{j+k}$ follow from linearity and independence. Over $q$ periods the reported returns sum to the
true ones except at the edges, so their variance is $q\sigma^2$ up to terms of order one; in Lo's formula the denominator $q + 2\sum(q-k)\rho_k$ times
$\sigma^2\sum\theta_j^2$ is exactly the variance of the sum of $q$ reported returns, which is $q\sigma^2$ plus edge terms.
\end{proof}

With the profile $(0.5, 0.3, 0.2)$ the variance is multiplied by 0.38 and the volatility by 0.616; a true Sharpe ratio of 0.6 is reported as 0.97, with
first autocorrelations of 0.55 and 0.26. The chapter's illiquid book has exactly this profile, and its five-year history shows the effect:

\begin{center}
\small
\begin{tabular}{@{}lcc@{}}
\toprule
the smoothed book, 60 months & reported & true or corrected\\
\midrule
annual volatility & 7.4\% & 11.2\% (true), 12.7\% (unsmoothed)\\
Sharpe ratio & 0.99 & 0.66 (true), 0.61 (Lo), 0.59 (unsmoothed)\\
standard error of the Sharpe ratio & 0.46 (iid) & 0.63 (HAC)\\
first two autocorrelations & 0.64, 0.25 & \\
smoothing profile & & $(0.43, 0.36, 0.21)$ estimated, $(0.5, 0.3, 0.2)$ true\\
\bottomrule
\end{tabular}
\end{center}

Three corrections agree: Lo's \omterm{def:rs:performance-measurement:annualisation}{annualisation} with the sample autocorrelations, the unsmoothed series (the moving average inverted with the profile estimated
from the first two autocorrelations, \cref{lst:rs:performance-measurement:unsmooth}), and the HAC standard error, which says that the reported 0.99 is also
less precise than it looks. Lo found hedge-fund Sharpe ratios overstated by up to 65\% through serial correlation; Getmansky, Lo and Makarov traced it to
illiquidity.

\section{Drawdowns}

\begin{definition}[Maximum drawdown, drawdown duration, Calmar ratio]\label{def:rs:performance-measurement:drawdown}
The \emph{maximum drawdown}\index{maximum drawdown} of a history is its largest fall of wealth from a running peak, as a fraction of the peak (drawdown: Book~2,
chapter~29). The \emph{drawdown duration}\index{drawdown duration} of a spell is the time from the peak to the recovery of that peak, or to the end of the history
if it has not recovered. The \emph{Calmar ratio}\index{Calmar ratio} is the annualised return over the absolute maximum drawdown, traditionally on three years.
\end{definition}

Drawdowns are what investors live through, and what ends strategies; they are also the least stable of the tear sheet's numbers, because one path yields one
maximum. For a Brownian motion with positive drift the expected \omterm{def:rs:performance-measurement:drawdown}{maximum drawdown} grows only logarithmically with the horizon (Magdon-Ismail, Atiya, Pratap and
Abu-Mostafa); a single history's maximum is one draw from a wide distribution. The trend stream's 33.6\% in five years is a bad draw: its model gives a 30\%
drawdown within five years with probability 14\%. The short-volatility stream's 14.8\% is a good one: 41\%. Over 200 simulated years the two reach \omterm{def:rs:performance-measurement:drawdown}{maximum
drawdowns} of 37\% and 72\%, with Sharpe ratios of 0.83 and 0.62 (\cref{fig:rs:performance-measurement:paths,fig:rs:performance-measurement:ddprob}). The \omterm{def:rs:performance-measurement:drawdown}{Calmar
ratios} of the displayed histories, 0.43 and 0.78, rank them the wrong way round.

\begin{omfigure}
\begin{tikzpicture}
\begin{axis}[name=top,width=12cm,height=5.2cm,ylabel={\small wealth},tick label style={font=\footnotesize},xmin=0,xmax=5,ymin=0.85,ymax=2.9,
  grid=major,grid style={gray!25},xticklabels={},legend columns=4,
  legend style={font=\footnotesize,at={(0.5,1.03)},anchor=south,draw=none,fill=none,/tikz/every even column/.append style={column sep=6pt}},
  table/col sep=comma]
\addplot[black!60,thick] table[x=year,y=trend]{figdata/research/22-performance-measurement/paths.csv};
\addplot[omThm,thick] table[x=year,y=shortvol]{figdata/research/22-performance-measurement/paths.csv};
\addplot[omDef,thick] table[x=year,y=smoothed]{figdata/research/22-performance-measurement/paths.csv};
\addplot[green!45!black,thick,dashed] table[x=year,y=mm]{figdata/research/22-performance-measurement/paths.csv};
\legend{trend,short volatility,smoothed book,market maker}
\end{axis}
\begin{axis}[at={(top.south west)},anchor=north west,yshift=-0.2cm,width=12cm,height=4.2cm,xlabel={\small years},ylabel={\small drawdown},
  tick label style={font=\footnotesize},xmin=0,xmax=5,ymin=-0.36,ymax=0.02,grid=major,grid style={gray!25},
  yticklabel style={/pgf/number format/fixed,/pgf/number format/precision=2},scaled y ticks=false,table/col sep=comma]
\addplot[black!60,thick] table[x=year,y=dd_trend]{figdata/research/22-performance-measurement/paths.csv};
\addplot[omThm,thick] table[x=year,y=dd_shortvol]{figdata/research/22-performance-measurement/paths.csv};
\addplot[omDef,thick] table[x=year,y=dd_smoothed]{figdata/research/22-performance-measurement/paths.csv};
\addplot[green!45!black,thick,dashed] table[x=year,y=dd_mm]{figdata/research/22-performance-measurement/paths.csv};
\end{axis}
\end{tikzpicture}

\omcaption{The four displayed five-year histories: wealth (top) and drawdown from the running peak (bottom). The smoothed book is monthly and drawn as steps
of a month. Data: \texttt{rs\_perf.histories}.}\label{fig:rs:performance-measurement:paths}
\end{omfigure}

\begin{omfigure}
\begin{tikzpicture}
\begin{axis}[width=10.5cm,height=5.4cm,xlabel={horizon (years)},ylabel={P(drawdown $\ge$ 30\%)},tick label style={font=\small},label style={font=\small},
  xmin=0.5,xmax=10.5,ymin=0,ymax=0.7,ybar,bar width=8pt,grid=major,grid style={gray!25},xtick={1,2,3,4,5,6,7,8,9,10},legend columns=2,
  yticklabel style={/pgf/number format/fixed,/pgf/number format/precision=1},scaled y ticks=false,
  legend style={font=\footnotesize,at={(0.5,-0.3)},anchor=north,draw=none,fill=none,/tikz/every even column/.append style={column sep=8pt}},
  table/col sep=comma]
\addplot[fill=gray!45,draw=gray] table[x=years,y=trend]{figdata/research/22-performance-measurement/ddprob.csv};
\addplot[fill=omThm!55,draw=omThm] table[x=years,y=shortvol]{figdata/research/22-performance-measurement/ddprob.csv};
\legend{trend,short volatility}
\end{axis}
\end{tikzpicture}

\omcaption{The probability of a drawdown of 30\% or more within a horizon, over 4\,000 simulated paths of each strategy's model (long-run volatilities of 15\% and
17\%). Data: \texttt{rs\_perf.dd\_probability}.}\label{fig:rs:performance-measurement:ddprob}
\end{omfigure}

\section{Turnover, hit rate and holding period}

\begin{definition}[Hit rate, profit factor]\label{def:rs:performance-measurement:hit}
The \emph{hit rate}\index{hit rate} of a return stream is the share of its non-zero periods with a positive return. Its \emph{profit factor}\index{profit factor}
is the sum of its gains over the sum of its losses.
\end{definition}

A \omterm{def:rs:performance-measurement:hit}{hit rate} says how often, not how much: the short-volatility stream wins on 76\% of days and the trend stream on 53\%, and the second has the better tail. A
\omterm{def:rs:performance-measurement:hit}{profit factor} above one is a positive mean and nothing more; it equals the \omterm{def:rs:performance-measurement:sortino}{Omega ratio} at zero (next section), since both are the mean gain over the mean loss.
The market maker's 65\% of winning days and \omterm{def:rs:performance-measurement:hit}{profit factor} of 2.07 are those of a business collecting a spread with rare adverse days; its Sharpe ratio of 4.04
(standard error 0.45, HAC 0.52) is real and says nothing about capacity (chapter~28).

Turnover (chapter~16) sets costs, and with the gross exposure it sets the average holding period: a book of gross exposure $G$ that trades $T$ of capital one
way per period replaces itself every $G/T$ periods. Chapter~16's momentum book, rebuilt monthly on the 500 most liquid names, trades 1.86\% of capital a day and
holds a position for 48 trading days on average; its Sharpe ratio after costs is 0.17 with a standard error of 0.32 over ten years, and its longest drawdown lasted 2\,079 of its
2\,520 days. A strategy's holding period should match its signal's half-life (chapter~13); a holding period much shorter than the half-life is
turnover paid for noise.

\section{The shape of outcomes}

\begin{definition}[Sortino ratio, Omega ratio]\label{def:rs:performance-measurement:sortino}
The \emph{Sortino ratio}\index{Sortino ratio} is the mean return in excess of a minimum acceptable return, divided by the downside deviation, the root mean square
of the shortfalls below it (Sortino and Price). The \emph{Omega ratio}\index{Omega ratio} at a threshold $L$ is $\mathbb{E}[(r-L)^+]\,/\,\mathbb{E}[(L-r)^+]$,
the expected gain above the threshold over the expected loss below it (Keating and Shadwick); it equals one at the mean.
\end{definition}

Downside measures were designed to stop penalising strategies for their upside. They cannot see a downside the sample does not contain. In the displayed history
the short-volatility stream's daily returns have a skewness of $+5.38$ and a kurtosis of 102.5; over 100 simulated years the skewness is $-1.23$ and the kurtosis
121. The five years held rebounds, not crashes, and its \omterm{def:rs:performance-measurement:sortino}{Sortino ratio} (1.87) and \omterm{def:rs:performance-measurement:sortino}{Omega ratio} (1.49) beat the trend stream's (1.48 and 1.17). Selling options
raises the Sharpe ratio and the downside ratios of a history that has not yet met the options' payoff; Goetzmann, Ingersoll, Spiegel and Welch showed how any
of the popular measures can be gamed this way, and characterised a measure that cannot be: an average of a power utility over the return history.

The remedy is to measure the strategy, not the history: to know what the strategy is short of, to simulate or stress it where the history is silent (Book~6's
stress tests), and to read shape statistics as descriptions of a sample.

\section{Benchmark-relative measures}

\begin{definition}[Market beta]\label{def:rs:performance-measurement:beta}
A strategy's \emph{market beta}\index{market beta} is the slope of its returns on the market's, $\beta = \operatorname{Cov}(r, r_m)/\operatorname{Var}(r_m)$,
estimated by regression with an intercept, the alpha (Book~1, chapter~1); standard errors are HAC.
\end{definition}

Regressed on the index it sells puts on, the short-volatility stream's displayed history has a beta of 0.45 (standard error 0.09) and an annual alpha of 5.6\%
($t = 1.38$). A beta is a linear summary of a nonlinear exposure. Estimated separately on the index's down days and up days, the history's betas are 0.59 and
0.54; over 100 simulated years they are 0.98 and 0.71, and the monthly returns trace the payoff of a short put (\cref{fig:rs:performance-measurement:scatter}):
flat above, steep below. The trend stream's beta is 0.00 (0.03), with an alpha $t$-statistic of 2.15 over five years; the market maker's 0.04 (0.01), $t =
9.1$. Against its benchmark a manager reports the information ratio and the tracking error (Book~1, chapter~3); the same caution applies.

\begin{omfigure}
\begin{tikzpicture}
\begin{axis}[width=9cm,height=6cm,xlabel={index monthly return},ylabel={short-volatility monthly return},tick label style={font=\small},label style={font=\small},
  xmin=-0.19,xmax=0.14,ymin=-0.45,ymax=0.2,grid=major,grid style={gray!25},
  xticklabel style={/pgf/number format/fixed,/pgf/number format/precision=2},yticklabel style={/pgf/number format/fixed,/pgf/number format/precision=1},
  scaled ticks=false,table/col sep=comma]
\addplot[only marks,mark=*,mark size=0.7pt,omThm] table[x=index,y=shortvol]{figdata/research/22-performance-measurement/scatter.csv};
\end{axis}
\end{tikzpicture}

\omcaption{Monthly returns of the short-volatility strategy against the index, over 100 simulated years: a short put's payoff, with daily betas of 0.98 on
down days and 0.71 on up days. Data: \texttt{rs\_perf.monthly\_scatter}.}\label{fig:rs:performance-measurement:scatter}
\end{omfigure}

\section{Tutorial: four streams and a book}\label{tut:rs:performance-measurement:tut}

\begin{tutorial}
\textbf{Goal.} Compute \texttt{firm.perf}'s tear sheet for the four synthetic streams and chapter~16's momentum book, and read what each number hides.
\textbf{End state:} the tear-sheet table below, \cref{fig:rs:performance-measurement:paths,fig:rs:performance-measurement:ddprob,fig:rs:performance-measurement:scatter}.
\begin{tutsteps}
\item \textbf{Lo's \omterm{def:rs:performance-measurement:annualisation}{annualisation}}: the per-period Sharpe ratio times $q/\sqrt{q + 2\sum(q-k)\rho_k}$.
\omcode{code/firm/perf/firm_perf.py}{63}{72}{The autocorrelation-adjusted Sharpe ratio.}\label{lst:rs:performance-measurement:lo}
\item \textbf{Drawdowns and their spells.}
\omcode{code/firm/perf/firm_perf.py}{79}{105}{Drawdown, maximum drawdown and drawdown spells.}\label{lst:rs:performance-measurement:dd}
\item \textbf{Unsmoothing}: invert the moving average with the estimated profile.
\omcode{code/firm/perf/firm_perf.py}{175}{183}{Recovering true returns from smoothed ones.}\label{lst:rs:performance-measurement:unsmooth}
\item \textbf{Run} \texttt{rs\_perf.histories()}, \texttt{stream\_results()}, \texttt{momentum()} and \texttt{tear\_sheet} on each, then \texttt{smoothed\_report()},
\texttt{dd\_probability()} and \texttt{fig\_perf.py}. The tear sheets, for five years of daily returns (the smoothed book: 60 months) and the momentum book's
ten years, with betas against the short-volatility simulation's index and, for momentum, the synthetic market:
\end{tutsteps}
\begin{center}
\small
\begin{tabular}{@{}lrrrrr@{}}
\toprule
 & trend & short vol. & smoothed & market maker & momentum\\
\midrule
annual return & 14.4\% & 11.6\% & 7.2\% & 22.8\% & 1.1\%\\
annual volatility & 14.6\% & 10.6\% & 7.4\% & 5.1\% & 9.1\%\\
Sharpe ratio & 1.00 & 1.09 & 0.99 & 4.04 & 0.17\\
\quad standard error, iid & 0.45 & 0.45 & 0.46 & 0.45 & 0.32\\
\quad standard error, HAC & 0.46 & 0.45 & 0.63 & 0.52 & 0.31\\
\omterm{def:rs:performance-measurement:drawdown}{maximum drawdown} & $-33.6$\% & $-14.8$\% & $-11.0$\% & $-4.2$\% & $-34.0$\%\\
longest drawdown & 791 d & 48 d & 22 m & 72 d & 2\,079 d\\
Calmar & 0.43 & 0.78 & 0.66 & 5.39 & 0.03\\
Sortino & 1.48 & 1.87 & 2.00 & 6.31 & 0.24\\
Omega at zero (\omterm{def:rs:performance-measurement:hit}{profit factor}) & 1.17 & 1.49 & 2.07 & 2.07 & 1.03\\
skewness & $-0.04$ & 5.38 & 0.41 & $-0.66$ & 0.00\\
kurtosis & 3.0 & 102.5 & 2.8 & 9.9 & 3.4\\
\omterm{def:rs:performance-measurement:hit}{hit rate} & 53\% & 76\% & 58\% & 65\% & 51\%\\
\omterm{def:rs:performance-measurement:beta}{market beta} (HAC s.e.) & 0.00 (0.03) & 0.45 (0.09) & & 0.04 (0.01) & 0.00 (0.01)\\
\bottomrule
\end{tabular}
\end{center}

\textbf{What to change next.} Double the short-volatility notional and compare the displayed Sharpe ratio with the probability of a 30\% drawdown; smooth the trend stream with the profile $(0.5, 0.3, 0.2)$ and see its Sharpe ratio rise by about 62\%.
\end{tutorial}



\section{Build: the performance module}\label{bld:rs:performance-measurement:perf}

\begin{build}
\bfield{Purpose} One function reads any \texttt{BacktestResult} (or return stream) and reports the numbers of this chapter with their uncertainty, so that no
strategy is judged on its Sharpe ratio alone.
\bfield{Interface} \texttt{from\_returns(r, periods)}, \texttt{sharpe}, \texttt{lo\_sharpe(r, q, lags)}, \texttt{drawdown}, \texttt{max\_drawdown},
\texttt{drawdown\_spells}, \texttt{longest\_drawdown}, \texttt{calmar}, \texttt{sortino(r, periods, mar)}, \texttt{omega(r, threshold)}, \texttt{hit\_rate},
\texttt{profit\_factor}, \texttt{holding\_period(res)}, \texttt{alpha\_beta(r, bench, periods, lags)}, \texttt{smoothing\_profile(r, k)},
\texttt{unsmooth(r, theta)}, \texttt{tear\_sheet(res, bench, name, periods)}.
\bfield{Rules} Every Sharpe ratio is printed with its standard error, iid and HAC; monthly or smoothed returns are annualised with their autocorrelations;
drawdown probabilities come from a model of the strategy, not from its history's maximum; shape statistics are labelled as the sample's.
\bfield{Acceptance tests} \texttt{code/firm/perf/tests/}: drawdowns, spells, Calmar, Sortino, Omega, \omterm{def:rs:performance-measurement:hit}{hit rate} and \omterm{def:rs:performance-measurement:hit}{profit factor} by hand; Lo's ratio equal to
$\sqrt q\,\mathrm{SR}_1$ on independent returns and correcting a smoothed stream; the smoothing profile recovered within 0.05 and the true returns recovered
exactly after the start-up; a known beta and alpha; a holding period from known turnover.
\bfield{Stretch} Deflated Sharpe ratios from the \omterm{def:rs:the-research-process:log}{research log}'s \omterm{def:rs:the-research-process:log}{trial count} (chapter~20); rolling tear sheets; the manipulation-proof measure.
\end{build}

\begin{omsources}
\item A.~W.~Lo, ``The statistics of Sharpe ratios'', \emph{Financial Analysts Journal} 58(4), 2002.
\item M.~Getmansky, A.~W.~Lo and I.~Makarov, ``An econometric model of serial correlation and illiquidity in hedge fund returns'', \emph{Journal of Financial
Economics} 74(3), 2004.
\item M.~Magdon-Ismail, A.~F.~Atiya, A.~Pratap and Y.~S.~Abu-Mostafa, ``On the maximum drawdown of a Brownian motion'', \emph{Journal of Applied Probability}
41(1), 2004.
\item F.~Sortino and L.~Price, ``Performance measurement in a downside risk framework'', \emph{Journal of Investing} 3(3), 1994; PerformanceAnalytics (R),
documentation of \texttt{SortinoRatio}, \texttt{CalmarRatio} and \texttt{Omega}.
\item W.~F.~Sharpe, ``Mutual fund performance'', \emph{Journal of Business} 39(S1), 1966; ``The Sharpe ratio'', \emph{Journal of Portfolio Management}
21(1), 1994.
\item W.~Goetzmann, J.~Ingersoll, M.~Spiegel and I.~Welch, ``Portfolio performance manipulation and manipulation-proof performance measures'', \emph{Review of
Financial Studies} 20(5), 2007.
\end{omsources}

\section{Exercises}

\begin{exercise}[$\star$]\label{exo:rs:performance-measurement:1}
A strategy has a Sharpe ratio of 1.0 over five years of daily returns. What is its standard error, and its 95\% interval?
\end{exercise}

\begin{exercise}[$\star$]\label{exo:rs:performance-measurement:2}
The trend stream returned 14.4\% a year with a \omterm{def:rs:performance-measurement:drawdown}{maximum drawdown} of 33.6\%. Compute its \omterm{def:rs:performance-measurement:drawdown}{Calmar ratio}. Why is the \omterm{def:rs:performance-measurement:drawdown}{Calmar ratio} of a single history unstable?
\end{exercise}

\begin{exercise}[$\star$]\label{exo:rs:performance-measurement:3}
Show that the \omterm{def:rs:performance-measurement:hit}{profit factor} equals the \omterm{def:rs:performance-measurement:sortino}{Omega ratio} at a threshold of zero.
\end{exercise}

\begin{exercise}[$\star\star$]\label{exo:rs:performance-measurement:4}
True returns with a Sharpe ratio of 0.6 are reported smoothed with the profile $(0.5, 0.3, 0.2)$. Give the factor on the variance, on the volatility, the reported
Sharpe ratio, and the first two autocorrelations.
\end{exercise}

\begin{exercise}[$\star\star$]\label{exo:rs:performance-measurement:5}
With the autocorrelations of exercise~4 (and zero beyond lag 2), compute Lo's \omterm{def:rs:performance-measurement:annualisation}{annualisation} factor for monthly returns and compare it with $\sqrt{12}$. By how
much does the naive \omterm{def:rs:performance-measurement:annualisation}{annualisation} overstate the Sharpe ratio?
\end{exercise}

\begin{exercise}[$\star\star$]\label{exo:rs:performance-measurement:6}
Why does a beta of 0.45 understate the short-volatility strategy's exposure to a crash? Use the down- and up-day betas.
\end{exercise}

\begin{exercise}[$\star\star\star$]\label{exo:rs:performance-measurement:7}
\emph{Coding.} With \texttt{rs\_perf.dd\_probability}, give the probability of a 30\% drawdown within one year and within ten for each strategy, and explain why
the short-volatility strategy's probability starts so much higher.
\end{exercise}

\begin{exercise}[$\star\star\star$]\label{exo:rs:performance-measurement:8}
\emph{Find the flaw.} ``Our fund's \omterm{def:rs:performance-measurement:sortino}{Sortino ratio} is 1.9 and it has never had a losing quarter in five years, so its downside risk is lower than that of the
trend fund next door, whose \omterm{def:rs:performance-measurement:sortino}{Sortino ratio} is 1.5 and which lost a third of its value.''
\end{exercise}

\section{Problem: Two Sharpe Ratios of One}

\begin{problem}[{Weekend problem --- two Sharpe ratios of one}]\label{pb:rs:performance-measurement:1}
The chapter's four streams and chapter~16's momentum book, through \texttt{firm.perf}.

\medskip\noindent\textbf{Part I --- The Sharpe ratio.}
\begin{enumerate}
\item What are the Sharpe ratios of the trend and short-volatility histories, and their standard errors?
\item Why do four of the five streams have nearly the same standard error?
\item What does the HAC standard error say about the smoothed book?
\item What is the market maker's Sharpe ratio and its standard error, and what does it not tell you?
\end{enumerate}

\medskip\noindent\textbf{Part II --- Smoothing.}
\begin{enumerate}[resume]
\item What are the smoothed book's reported and true volatilities?
\item What are its reported, Lo-adjusted, unsmoothed and true Sharpe ratios?
\item Which smoothing profile is estimated, from which autocorrelations?
\item Why is the naive \omterm{def:rs:performance-measurement:annualisation}{annualisation} wrong for this book?
\end{enumerate}

\medskip\noindent\textbf{Part III --- Drawdowns and shape.}
\begin{enumerate}[resume]
\item Give each displayed history's \omterm{def:rs:performance-measurement:drawdown}{maximum drawdown} and longest drawdown.
\item Which strategy has the better Calmar, Sortino and \omterm{def:rs:performance-measurement:sortino}{Omega ratios} in the displayed histories, and which has the worse tail?
\item What are the short-volatility stream's skewness and kurtosis in the history and over 100 years?
\item What share of its five-year histories have no losing quarter?
\item What are the momentum book's holding period and turnover?
\end{enumerate}

\medskip\noindent\textbf{Part IV --- The verdict.}
\begin{enumerate}[resume]
\item State the \emph{named result}: the probability of a 30\% drawdown within five years for each strategy, and the Sharpe ratio's standard error under the
smoothed book's autocorrelation.
\item What do 200 simulated years give for each strategy's Sharpe ratio and \omterm{def:rs:performance-measurement:drawdown}{maximum drawdown}?
\item What are the short-volatility strategy's beta, down-day and up-day betas?
\item Which of the two would you allocate to, and how would you size it?
\item What should every tear sheet print next to a Sharpe ratio?
\item How could a manager game the measures of this chapter?
\item In one sentence: what does a Sharpe ratio of one over five years tell you?
\end{enumerate}
\end{problem}

\section{Interview questions}

\begin{interviewq}[$\star$ \iqroles{researcher, trader}]\label{iq:rs:performance-measurement:1}
A strategy shows a Sharpe ratio of 1 over three years of daily data. How confident are you that it is positive?
\end{interviewq}

\begin{interviewq}[$\star\star$ \iqroles{researcher, risk}]\label{iq:rs:performance-measurement:2}
How can a strategy raise its Sharpe ratio without adding skill? How would you detect it?
\end{interviewq}

\begin{interviewq}[$\star\star$ \iqroles{researcher}]\label{iq:rs:performance-measurement:3}
A fund's monthly returns have a first autocorrelation of 0.6. What does that do to its Sharpe ratio, and how do you correct it?
\end{interviewq}

\begin{interviewq}[$\star\star$ \iqroles{risk}]\label{iq:rs:performance-measurement:4}
How does the expected \omterm{def:rs:performance-measurement:drawdown}{maximum drawdown} of a strategy grow with the horizon? Why is a \omterm{def:rs:vectorised-backtests:backtest}{backtest}'s \omterm{def:rs:performance-measurement:drawdown}{maximum drawdown} a poor risk estimate?
\end{interviewq}

\begin{interviewq}[$\star\star$ \iqroles{trader}]\label{iq:rs:performance-measurement:5}
A market maker wins on 65\% of days with a Sharpe ratio of 4. What would you still want to know?
\end{interviewq}

\begin{interviewq}[$\star\star\star$ \iqroles{researcher}]\label{iq:rs:performance-measurement:6}
Derive the standard error of the Sharpe ratio for independent normal returns.
\end{interviewq}
